\documentclass[11pt]{article}
\usepackage[margin=1in]{geometry}
\usepackage{amsmath}
\usepackage{natbib}
\usepackage[hidelinks]{hyperref}

\title{Noise Geometry and Escape Dynamics in SGD and SGLD on Nonsmooth
Nonconvex Landscapes: A Controlled Empirical Study}
\author{Research draft for internal review}
\date{}

\begin{document}
\maketitle

\section{Introduction}

The same perturbation that helps an optimizer move away from a saddle may
hinder its subsequent progress. A large update can cross the boundary of a
saddle neighborhood without moving in a direction of descent or improving
the objective. Conversely, leaving that neighborhood with a lower training
loss need not improve performance on unseen data. These distinctions matter
when comparing stochastic gradient descent (SGD), whose randomness comes from
sampling training components, with stochastic gradient Langevin dynamics
(SGLD), which adds an explicit, typically isotropic, Gaussian perturbation
\citep{raginsky2017,zhu2019}.

Nonsmoothness complicates the measurement. The ordinary Hessian may be
undefined at a kink, and a subgradient sampled off the kink need not be small
even arbitrarily close to a critical point. Active-manifold analyses
distinguish the tangential geometry of a nonsmooth saddle from the behavior
across the kink \citep{bianchi2023}. This distinction motivates our
synthetic benchmarks. For models whose active geometry is not established, we
plan to use explicitly labeled diagnostic proxies rather than treat their
outputs as certificates of a strict saddle.

We ask when the \emph{structure, intensity, and timing} of stochastic
perturbations improve productive escape, subsequent optimization, and
out-of-sample quality at equal component-oracle cost. We plan three elements:
first, define and record geometric exit, movement along an unstable
direction, and exit with a prespecified objective decrease as distinct
events; second, compare intrinsic mini-batch noise and isotropic perturbations
under controls for total noise scale and the cost of its calibration; third,
test whether early versus sustained noise, ambient dimension, and changes in
batch size alter the relationship between escape and final quality. Analytical
landscapes are intended to isolate mechanisms, while statistical tasks are
needed to evaluate out-of-sample quality. This is a proposed empirical design,
not a claim of universal superiority or a report of experimental findings.

\section{Related work}

\paragraph{Escaping saddles in smooth and nonsmooth optimization.}
For smooth objectives, perturbed gradient descent can find approximate
second-order stationary points under Lipschitz-gradient and
Lipschitz-Hessian assumptions \citep{jin2017}. \citet{daneshmand2018} study
the second moment of stochastic gradients along a negative-curvature
direction, while \citet{fang2019} use a dispersiveness condition on the
gradient noise. These conditions motivate directional measurements, but do
not identify the probability of our finite-budget, loss-qualified exit
event. In nonsmooth optimization, \citet{bianchi2023} establish
nonconvergence to active strict saddles under assumptions on the geometry,
steps, and perturbations. \citet{davis2021} instead study finite-time escape
from strict saddles of the Moreau envelope using inexact proximal oracles.
Our geometric classification of iterates and finite-budget events do
not require interpreting an arbitrary ReLU iterate as a mathematically
certified active saddle.

\paragraph{Langevin methods, noise controls, and generalization.}
\citet{zhang2017} bound SGLD hitting times to selected sets; this differs
from observing a crossing of a local saddle neighborhood with a stipulated
decrease in loss. \citet{zhu2019} compare covariances with matched trace in
the context of escape from sharp minima, motivating a scale-matched control
without transferring their minimum-escape statistic to saddles.
\citet{mou2017} analyze generalization via stability and PAC-Bayesian
arguments, and \citet{zhang2020} alternate exploration and sampling in
cyclic stochastic-gradient MCMC. Proximal Langevin schemes provide another
way to handle nonsmooth potentials when the goal is sampling
\citep{durmus2016,renaud2025}. Our proposed comparison is an optimization
study: it evaluates exit and held-out quality separately, conditions
comparisons on declared oracle cost, and treats the allocation of a
matched amount of perturbation energy over time as a separate question.

\begin{thebibliography}{12}
\bibitem[Jin et al.(2017)]{jin2017}
C.~Jin et al. How to Escape Saddle Points Efficiently. 2017.
\url{https://www.alphaxiv.org/abs/1703.00887}
\bibitem[Daneshmand et al.(2018)]{daneshmand2018}
H.~Daneshmand et al. Escaping Saddles with Stochastic Gradients. 2018.
\url{https://www.alphaxiv.org/abs/1803.05999}
\bibitem[Fang et al.(2019)]{fang2019}
C.~Fang et al. Sharp Analysis for Nonconvex SGD Escaping from Saddle Points.
2019. \url{https://www.alphaxiv.org/abs/1902.00247}
\bibitem[Bianchi et al.(2023)]{bianchi2023}
P.~Bianchi, W.~Hachem, and S.~Schechtman. Stochastic Subgradient Descent
Escapes Active Strict Saddles on Weakly Convex Functions. 2023 preprint.
\url{https://www.alphaxiv.org/abs/2108.02072}
\bibitem[Davis et al.(2021)]{davis2021}
D.~Davis, M.~D\'iaz, and D.~Drusvyatskiy. Escaping Strict Saddle Points of
the Moreau Envelope in Nonsmooth Optimization. 2021 preprint.
\url{https://www.alphaxiv.org/abs/2106.09815}
\bibitem[Zhang et al.(2017)]{zhang2017}
Y.~Zhang, P.~Liang, and M.~Charikar. A Hitting Time Analysis of Stochastic
Gradient Langevin Dynamics. 2017. \url{https://www.alphaxiv.org/abs/1702.05575}
\bibitem[Raginsky et al.(2017)]{raginsky2017}
M.~Raginsky, A.~Rakhlin, and M.~Telgarsky. Non-Convex Learning via
Stochastic Gradient Langevin Dynamics: A Nonasymptotic Analysis. 2017.
\url{https://www.alphaxiv.org/abs/1702.03849}
\bibitem[Mou et al.(2017)]{mou2017}
W.~Mou et al. Generalization Bounds of SGLD for Non-convex Learning:
Two Theoretical Viewpoints. 2017 preprint.
\url{https://www.alphaxiv.org/abs/1707.05947}
\bibitem[Zhu et al.(2019)]{zhu2019}
Z.~Zhu et al. The Anisotropic Noise in Stochastic Gradient Descent:
Its Behavior of Escaping from Sharp Minima and Regularization Effects. 2019.
\url{https://www.alphaxiv.org/abs/1803.00195}
\bibitem[Zhang et al.(2020)]{zhang2020}
R.~Zhang et al. Cyclical Stochastic Gradient MCMC for Bayesian Deep
Learning. 2020. \url{https://www.alphaxiv.org/abs/1902.03932}
\bibitem[Durmus et al.(2016)]{durmus2016}
A.~Durmus et al. Efficient Bayesian Computation by Proximal Markov Chain
Monte Carlo: When Langevin Meets Moreau. 2016 preprint.
\url{https://www.alphaxiv.org/abs/1612.07471}
\bibitem[Renaud et al.(2025)]{renaud2025}
M.~Renaud et al. From Stability of Langevin Diffusion to Convergence of
Proximal MCMC for Non-log-concave Sampling. 2025 preprint.
\url{https://www.alphaxiv.org/abs/2505.14177}
\end{thebibliography}
\end{document}
